% !TEX program = lualatex
% !TeX encoding = UTF-8
\documentclass[11pt,a4paper]{article}

% ============================================================
% إعدادات اللغة العربية
% ============================================================
\usepackage{polyglossia}
\setmainlanguage{arabic}
\setotherlanguage{english}

% ========= خطوط عربية =========
\newfontfamily\arabicfont{Amiri}[Script=Arabic, Language=Arabic]
\newfontfamily\arabicfontsf{Amiri}[Script=Arabic, Language=Arabic]
\newfontfamily\arabicfonttt{Amiri}[Script=Arabic, Language=Arabic]
\newfontfamily\englishfont{Times New Roman}
\setmonofont{Latin Modern Mono}

% ========= إزالة تحذيرات hyperref =========
\makeatletter
\let\@ensure@LTR\relax
\providecommand{\babelsublr}[1]{#1}
\makeatother

% ========= حزم =========
\usepackage{amsmath,amssymb,amsthm,mathtools}
\usepackage{geometry}
\geometry{left=1.5cm,right=1.5cm,top=1.5cm,bottom=2.0cm}
\usepackage{booktabs}
\usepackage{longtable}
\usepackage{array}
\usepackage{hyperref}
\usepackage{url}
\usepackage{xcolor}
\usepackage{fancyhdr}
\usepackage{enumitem}
\usepackage{makecell}
\usepackage{float}

% ========= ألوان =========
\definecolor{darkblue}{RGB}{0,0,139}
\definecolor{skyblue}{RGB}{0,102,204}
\definecolor{gold}{RGB}{255,215,0}

% ========= إعدادات الوصلات التشعبية =========
\hypersetup{colorlinks=true,linkcolor=blue,citecolor=blue,urlcolor=blue}

% ========= ثوابت YUCT =========
\newcommand{\REL}{\mathcal{K}}
\newcommand{\ABS}{K_{\mathrm{eff}}}
\newcommand{\kappac}{\kappa_c}
\newcommand{\betaf}{\beta}
\newcommand{\Sodd}{S_{\mathrm{odd}}}
\newcommand{\Seven}{S_{\mathrm{even}}}

\begin{document}
	
	\begin{titlepage}
		\centering
		\vspace*{2cm}
		{\Huge\bfseries الملحق BCLM. لاغرانجية التنسيق البيولوجي\par}
		\vspace{0.3cm}
		{\Large\bfseries (الإصدار 6.1 – تعريفات صارمة، معايرة شفافة، وتنبؤات مصاغة بحذر)\par}
		\vspace{1.5cm}
		{\large A.\,V.\,Yakushev\par}
		{\normalsize YUCT Research, \url{https://yuct.org/}\par}
		
		\vspace{1.0cm}
		{\normalsize \textbf{DOI:} \url{https://doi.org/10.5281/zenodo.18444598}\par}
		\vspace{1.0cm}
		{\normalsize أيار 2026\par}
		\vfill
		\begin{abstract}
			\noindent
			يقدم هذا المستند لاغرانجية التنسيق البيولوجي (BCLM) كوحدة متسقة تعتمد على البيانات من نظرية التنسيق الموحد لياكوشيف (YUCT). 
			يُعطى كل كائن بيولوجي كفاءة تنسيق نسبية \(\REL\)، تُعرف بشكل موحد من خلال بروتوكول YPSDC كنسبة الضغط \(H(\mathcal{A})/H(\kappa)\) مطبعة على الحد الأقصى النظري. 
			نقدم جداول كاملة للكودونات، الأحماض الأمينية، الإنزيمات، كفاءات الإصلاح المرتبطة بالطفرات، الثوابت التباينية، ومعاملات الخلايا العصبية. 
			قاعدة توافق جبرية عالمية، مستمدة من ثوابت YUCT ومعايرة على بيانات ارتباط البروتين-بروتين، تحكم التفاعلات بين أي كيانين بيولوجيين. 
			تمت صياغة بروتوكول الحياة الطويلة كفرضية كمية قابلة للتكذيب؛ وتُشتق تنبؤاته من قانون الخطأ العالمي والبنية الضربية للأخطاء الخلوية.
			تم استعادة الوحدات المتقدمة الرئيسية (التنسيق الزمني، المناعة، السموم، الكائنات المتعددة، الشابرونات، التعديلات ما بعد الترجمة) والتعبير عنها في إطار \(\REL\) الموحد.
			تمت مناقشة جميع التعريفات والمعايرات والقيود لضمان الدقة العلمية.
		\end{abstract}
	\end{titlepage}
	
	\tableofcontents
	\newpage
	
	% ============================================================
	% 1. مقدمة: من PLCM إلى علم الأحياء
	% ============================================================
	\section{مقدمة: من PLCM إلى علم الأحياء}
	تصف نظرية التنسيق YUCT أي نظام معقد من خلال المعامل العالمي \(K_{\mathrm{eff}}\) — كفاءة التنسيق. بالنسبة للعناصر الكيميائية، تم تنفيذ هذه الفكرة في الملحق PLCM، حيث تم تعيين \(K_{\mathrm{eff}}\) لكل عنصر وتم ترميز التفاعلات في مصفوفة اقتران \(\kappa_{ij}\). نقوم هنا بتوسيع نفس النهج ليشمل الأنظمة الحية.
	
	علم الأحياء منظم هرمياً: النيوكليوتيدات، الكودونات، الأحماض الأمينية، البروتينات، الشبكات الأيضية، الكائنات الحية، المجموعات السكانية. في كل مستوى، يمكن تعريف \(K_{\mathrm{eff}}\) عبر بروتوكول YPSDC: مؤشر قصير \(\kappa\) (مثل كودون، جزيء ركيزة) يُفعِّل فعلاً معقداً \(\mathcal{A}\) (مثل إدخال حمض أميني، تحويل تحفيزي) من قاموس موزع مسبقاً. تكون كفاءة التنسيق إذن
	\[
	K_{\mathrm{eff}} = \frac{H(\mathcal{A})}{H(\kappa)},
	\]
	حيث \(H\) تشير إلى إنتروبيا شانون. هذا التعريف التشغيلي لا يحتاج إلى معاملات حرة: يتم حساب كلا الإنتروبيين من توزيعات التردد التجريبية المتاحة للعموم.
	
	في هذا الملحق نستخدم \textbf{كفاءة التنسيق النسبية}
	\[
	\REL \equiv \frac{K_{\mathrm{eff}}}{K_{\mathrm{eff}}^{\max}} \in (0,1],
	\]
	حيث \(K_{\mathrm{eff}}^{\max}\) هو الحد الأقصى النظري للضغط الممكن لهذه الفئة من الكائنات (مثل 6 بت للكودون). جميع القيم لا أبعاد لها وقابلة للمقارنة مباشرة عبر المستويات.
	
	\section{الهيكل العام للاغرانجية البيولوجية}
	على غرار PLCM، فإن لاغرانجية التنسيق البيولوجي هي
	\begin{equation}
		\mathcal{L}_{\text{BCLM}} = \sum_{s=1}^{N_{\text{sectors}}} \lambda_s L_s(\Phi_s) + \sum_{s<r} \kappa_{sr} \operatorname{Tr}(\Psi_{sr} \cdot O_s \cdot O_r^\dagger),
		\label{eq:BCLM}
	\end{equation}
	حيث \(\Phi_s\) هي كميات قابلة للرصد في القطاع \(s\)، \(\kappa_{sr}\) معاملات الاقتران بين القطاعات، و \(\Psi_{sr}\) حقل التنسيق. المهمة الرئيسية هي حساب \(\REL\) لكل كائن بيولوجي وبناء مصفوفات التوافق \(C_{AB}\). التطابق مع لاغرانجية YUCT ذات 120 قطاعاً (الملحق A) هو:
	\begin{itemize}
		\item القطاع 40 (DNA) – الشفرة الوراثية، الكودونات؛
		\item القطاع 41 (RNA) – النسخ؛
		\item القطاع 42 (البروتينات) – الأحماض الأمينية، الطي؛
		\item القطاع 43 (التمثيل الغذائي) – الإنزيمات؛
		\item القطاعات 50–55 (علم الأعصاب) – الخلايا العصبية؛
		\item القطاعات 60–61 (التطور) – معدلات الطفرة، الانتواع.
	\end{itemize}
	
	% ============================================================
	% 2. تعريف صارم لـ \(\REL\)
	% ============================================================
	\section{تعريف \(\REL\): الأساس المعلوماتي}
	\label{sec:REL-def}
	
	\subsection{تعريف تشغيلي عام}
	يفصل بروتوكول YPSDC بين توزيع القاموس خارج الخط وإرسال مؤشر قصير على الخط. لأي عملية بيولوجية، نحدد:
	\begin{itemize}
		\item \(\mathcal{A}\) – مجموعة الأفعال الممكنة (النتائج) ذات الترددات النسبية المرصودة \(p_i\). إنتروبيا الفعل هي \(H(\mathcal{A}) = -\sum_i p_i \log_2 p_i\).
		\item \(\kappa\) – المؤشر الجزيئي الذي يختار مدخلاً معيناً في القاموس. إنتروبياه هي \(H(\kappa) = -\log_2 p_\kappa\)، حيث \(p_\kappa\) هو تردد ذلك المؤشر في المجموعة ذات الصلة (مثل استخدام الكودون، تركيز الأيض).
	\end{itemize}
	كفاءة التنسيق المطلقة هي \(\ABS = H(\mathcal{A})/H(\kappa)\). الكفاءة النسبية \(\REL = \ABS / \ABS^{\,\max}\) تطبع هذا إلى أقصى ضغط ممكن \(\ABS^{\,\max}\) يمكن أن يحمله الجزيء المؤشر فيزيائياً.
	
	\subsection{تطبيق على الكودونات}
	\textit{المؤشر:} ثلاثية كودون (64 احتمالاً)؛ ترددها \(p_\kappa\) مأخوذ من قاعدة بيانات استخدام الكودون البشري (Kazusa). \\
	\textit{الفعل:} إدخال حمض أميني في سلسلة بولي ببتيد نامية. توزيع الأحماض الأمينية في البروتيوم البشري (UniProt) يعطي \(p_i\) للأحماض الأمينية العشرين بالإضافة إلى التوقف. \\
	\textit{السعة القصوى:} \(\ABS^{\,\max}=6\) بت (ثلاثة أزواج قاعدية، كل منها على الأكثر 1 بت من المعلومات بعد التكرار). \\
	\textit{مثال:} الكودون CUG (Leu)، \(p_\kappa=0.0396\)، \(H(\kappa)\approx4.66\) بت؛ \(H(\mathcal{A})\approx4.19\) بت؛ \(\ABS=0.900\)، وبالتالي \(\REL=0.150\).
	
	\subsection{تطبيق على الأحماض الأمينية}
	\textit{المؤشر:} نوع الحمض الأميني؛ تردده هو كسر البروتيوم. \\
	\textit{الفعل:} مجموعة التشكلات المختلفة للسلسلة الجانبية (روتامرات) التي يمكن لتلك البقيمة الوصول إليها. إنتروبيا الفعل هي \(\log_2\)(عدد الروتامرات المرصودة في PDB). \\
	\textit{السعة القصوى:} \(\log_2 20 \approx 4.32\) بت (المعلومات اللازمة لتحديد واحد من عشرين بقيمة). \\
	\textit{مثال:} ألانين (روتامرات = 3، \(p=0.070\)) يعطي \(H(\mathcal{A})=1.585\) بت، \(H(\kappa)=3.84\) بت، \(\ABS=0.413\)، \(\REL=0.096\).
	
	\subsection{تطبيق على الإنزيمات}
	\textit{المؤشر:} جزيء الركيزة؛ تركيزه الفعال يعرّف \(H(\kappa) = -\log_2(K_M/[S])\)، مع \([S]=1\,\text{mM}\) كمرجع قياسي. \\
	\textit{الفعل:} التحويل التحفيزي؛ إنتروبيا الفعل هي لوغاريتم عدد معدلات التحول المميزة المرصودة في عائلة الإنزيم، مقرباً بـ \(\log_2(k_{\mathrm{cat}}\tau_0)\) مع \(\tau_0=1\,\text{s}\). \\
	\textit{السعة القصوى:} \(\log_2(10^4)\approx13.3\) بت (تقدير متحفظ لعدد حالات الانتقال المميزة التي يمكن للموقع النشط تمييزها). \\
	\textit{مثال:} كربونيك أنهيدراز II، \(k_{\mathrm{cat}}=1.0\times10^6\,\text{s}^{-1}\)، \(K_M=8.3\times10^{-3}\,\text{M}\)، يعطي \(H(\mathcal{A})=19.9\) بت، \(H(\kappa)=3.05\) بت، \(\ABS=6.52\)، \(\REL=0.490\).
	
	\subsection{الاتساق}
	يضمن هذا التعريف أن \(\REL\) كمية لا أبعاد لها، معلوماتية بحتة، قابلة للحساب فقط من الترددات التجريبية. لا توجد معاملات قابلة للتعديل تدخل في حسابها.
	
	% ============================================================
	% 3. جداول البيانات البيولوجية الكاملة
	% ============================================================
	\section{جداول البيانات البيولوجية الكاملة}
	\label{sec:tables}
	
	\subsection{كفاءة التنسيق النسبية للكودونات (البشر)}
	\label{sec:codons}
	يسرد الجدول~\ref{tab:codons} \(\REL\) لجميع الكودونات الـ 64، محسوبة كما هو موصوف في القسم~\ref{sec:REL-def}.
	
	\begin{longtable}{|c|c|c|c|}
		\caption{كفاءة التنسيق النسبية للكودون، \(\REL\)}\label{tab:codons}\\
		\hline
		\textbf{الكودون} & \textbf{AA} & \textbf{التردد (\%)} & \(\REL\) \\
		\hline
		\endfirsthead
		\hline
		\textbf{الكودون} & \textbf{AA} & \textbf{التردد (\%)} & \(\REL\) \\
		\hline
		\endhead
		\hline
		\multicolumn{4}{r}{يتبع في الصفحة التالية}\\
		\endfoot
		\hline
		\endlastfoot
		UUU & Phe & 1.76 & 0.120 \\
		UUC & Phe & 2.03 & 0.124 \\
		UUA & Leu & 0.76 & 0.099 \\
		UUG & Leu & 1.29 & 0.111 \\
		CUU & Leu & 1.30 & 0.112 \\
		CUC & Leu & 1.96 & 0.123 \\
		CUA & Leu & 0.72 & 0.098 \\
		CUG & Leu & 3.96 & 0.150 \\
		AUU & Ile & 1.60 & 0.117 \\
		AUC & Ile & 2.08 & 0.125 \\
		AUA & Ile & 0.75 & 0.099 \\
		AUG & Met & 2.20 & 0.127 \\
		GUU & Val & 1.10 & 0.107 \\
		GUC & Val & 1.45 & 0.114 \\
		GUA & Val & 0.71 & 0.098 \\
		GUG & Val & 2.81 & 0.136 \\
		UCU & Ser & 1.44 & 0.114 \\
		UCC & Ser & 1.76 & 0.120 \\
		UCA & Ser & 1.20 & 0.109 \\
		UCG & Ser & 0.44 & 0.089 \\
		CCU & Pro & 1.62 & 0.117 \\
		CCC & Pro & 1.98 & 0.123 \\
		CCA & Pro & 1.66 & 0.118 \\
		CCG & Pro & 0.68 & 0.097 \\
		ACU & Thr & 1.22 & 0.110 \\
		ACC & Thr & 1.89 & 0.122 \\
		ACA & Thr & 1.44 & 0.114 \\
		ACG & Thr & 0.59 & 0.094 \\
		GCU & Ala & 1.84 & 0.121 \\
		GCC & Ala & 2.76 & 0.135 \\
		GCA & Ala & 1.58 & 0.117 \\
		GCG & Ala & 0.74 & 0.099 \\
		UAU & Tyr & 1.22 & 0.110 \\
		UAC & Tyr & 1.53 & 0.116 \\
		UAA & STOP & 0.21 & 0.078 \\
		UAG & STOP & 0.08 & 0.068 \\
		CAU & His & 0.91 & 0.103 \\
		CAC & His & 1.18 & 0.109 \\
		CAA & Gln & 1.20 & 0.109 \\
		CAG & Gln & 2.40 & 0.130 \\
		AAU & Asn & 1.28 & 0.111 \\
		AAC & Asn & 1.92 & 0.122 \\
		AAA & Lys & 1.88 & 0.122 \\
		AAG & Lys & 2.08 & 0.125 \\
		GAU & Asp & 1.52 & 0.116 \\
		GAC & Asp & 1.96 & 0.123 \\
		GAA & Glu & 1.96 & 0.123 \\
		GAG & Glu & 2.68 & 0.134 \\
		UGU & Cys & 1.04 & 0.106 \\
		UGC & Cys & 1.28 & 0.111 \\
		UGA & STOP & 0.10 & 0.070 \\
		UGG & Trp & 1.32 & 0.112 \\
		CGU & Arg & 0.78 & 0.100 \\
		CGC & Arg & 1.16 & 0.108 \\
		CGA & Arg & 0.62 & 0.095 \\
		CGG & Arg & 1.04 & 0.106 \\
		AGU & Ser & 0.88 & 0.102 \\
		AGC & Ser & 1.44 & 0.114 \\
		AGA & Arg & 0.80 & 0.100 \\
		AGG & Arg & 1.08 & 0.107 \\
		\hline
	\end{longtable}
	
	\subsection{كفاءة التنسيق النسبية للأحماض الأمينية}
	\label{sec:aminoacids}
	يعطي الجدول~\ref{tab:amino} \(\REL\) للأحماض الأمينية العشرين القياسية.
	
	\paragraph*{مبرر إنتروبيا الفعل المختارة.}
	عدد الروتامرات المختلفة للسلسلة الجانبية هو التقدير الأكثر تحفظاً للفضاء التشكلي الذي يمكن لبقيمة الوصول إليه قبل الطي.
	يمكن استخراجه مباشرة من هياكل البروتين عالية الدقة (PDB)، مما يجعله تشغيلياً بالكامل وقابلاً للتكرار.
	اختيارات أخرى ممكنة – على سبيل المثال، يمكن للمرء أن يشمل ميل البقيمة للمشاركة في التحفيز، أو طاقة الذوبان الحرة، أو عدد أنماط الروابط الهيدروجينية المختلفة التي يمكن أن تشكلها.
	هذه البدائل ستؤدي إلى قيم عددية مختلفة قليلاً لـ \(\REL\) وقد تكون أكثر ملاءمة لبعض التطبيقات المتخصصة (مثل تقييم بقيعات الموقع النشط بشكل منفصل).
	في الوقت الحالي، نعتمد التعريف القائم على الروتامر بسبب بساطته ووضوحه؛ سيتم اختبار كفايته من خلال القوة التنبؤية لقواعد التوافق الناتجة (القسم~\ref{sec:compatibility}).
	
	\begin{longtable}{|c|c|c|c|}
		\caption{كفاءة التنسيق النسبية للأحماض الأمينية}\label{tab:amino}\\
		\hline
		\textbf{الحرف} & \textbf{الاسم} & \textbf{التردد (\%)} & \(\REL\) \\
		\hline
		\endfirsthead
		\hline
		\textbf{الحرف} & \textbf{الاسم} & \textbf{التردد (\%)} & \(\REL\) \\
		\hline
		\endhead
		\hline
		\multicolumn{4}{r}{يتبع في الصفحة التالية}\\
		\endfoot
		\hline
		\endlastfoot
		A & ألانين & 7.0 & 0.096 \\
		C & سيستئين & 2.3 & 0.082 \\
		D & أسبارتيت & 5.3 & 0.094 \\
		E & غلوتاميت & 6.3 & 0.098 \\
		F & فينيل ألانين & 4.0 & 0.090 \\
		G & غليسين & 6.7 & 0.055 \\
		H & هستيدين & 2.3 & 0.082 \\
		I & إيزوليوسين & 5.3 & 0.094 \\
		K & لايسين & 5.9 & 0.096 \\
		L & ليوسين & 9.9 & 0.109 \\
		M & ميثيونين & 2.3 & 0.082 \\
		N & أسباراجين & 4.1 & 0.091 \\
		P & برولين & 4.9 & 0.070 \\
		Q & غلوتامين & 4.1 & 0.091 \\
		R & أرجينين & 5.9 & 0.109 \\
		S & سيرين & 7.0 & 0.080 \\
		T & ثريونين & 5.9 & 0.084 \\
		V & فالين & 6.6 & 0.086 \\
		W & تريبتوفان & 1.1 & 0.061 \\
		Y & تيروسين & 3.2 & 0.080 \\
		\hline
	\end{longtable}
	
	\subsection{كفاءة التنسيق النسبية للإنزيمات}
	\label{sec:enzymes}
	يسرد الجدول~\ref{tab:enzymes} \(\REL\) لإنزيمات مختارة، محسوبة بتركيز ركيزة قياسي \([S]=1\,\text{mM}\).
	
	\begin{longtable}{|c|c|c|}
		\caption{كفاءة التنسيق النسبية للإنزيمات}\label{tab:enzymes}\\
		\hline
		\textbf{الإنزيم} & \(k_{\mathrm{cat}}/K_M\) (M\(^{-1}\)s\(^{-1}\)) & \(\REL\) \\
		\hline
		\endfirsthead
		\hline
		\textbf{الإنزيم} & \(k_{\mathrm{cat}}/K_M\) (M\(^{-1}\)s\(^{-1}\)) & \(\REL\) \\
		\hline
		\endhead
		\hline
		\multicolumn{3}{r}{يتبع في الصفحة التالية}\\
		\endfoot
		\hline
		\endlastfoot
		كربونيك أنهيدراز II & \(8.3\times10^7\) & 0.490 \\
		DNA بوليميراز I & \(5.0\times10^6\) & 0.378 \\
		كحول ديهيدروجينيز & \(1.2\times10^5\) & 0.284 \\
		تربسين & \(1.6\times10^3\) & 0.151 \\
		يورياز & \(1.0\times10^9\) & 0.522 \\
		كاتالاز & \(1.0\times10^7\) & 0.425 \\
		أسيتيل كولين إستيراز & \(2.0\times10^8\) & 0.502 \\
		ليزوزيم & \(5.0\times10^4\) & 0.235 \\
		ريبونوكلياز A & \(1.5\times10^3\) & 0.145 \\
		\hline
	\end{longtable}
	
	\subsection{كفاءة الإصلاح المرتبطة بالطفرات}
	\label{sec:mutations}
	يربط قانون الخطأ العالمي \(\varepsilon = \kappa_c \alpha K_{\mathrm{eff}}^{-\beta}\) (\(\beta=2/3\), \(\kappa_c=1/3\)) معدل الطفرة لكل جيل \(\mu\) بكفاءة التنسيق لآلية إصلاح DNA، \(K_{\mathrm{eff,repair}}\). لتجنب الدائرية، نعاير الثابت الخاص بالنظام \(\alpha_{\mathrm{repair}}\) باستخدام قدرة الإصلاح المقاسة \textit{in vitro} للخلايا البشرية (معبراً عنها بـ \(K_{\mathrm{eff,repair}}^{\mathrm{human}} \approx 0.62\) في وحدات نسبية، مستمدة من دقة إصلاح استئصال النوكليوتيدات). باستخدام \(\mu_{\mathrm{human}}=1.1\times10^{-8}\)، نحصل على \(\alpha_{\mathrm{repair}} \approx 0.01\). ثم يسرد الجدول~\ref{tab:mutations} \(K_{\mathrm{eff,repair}}\) لعدة مجموعات \textbf{متوقعة} من معدلات الطفرات تحت افتراض أن القانون صحيح. تعمل هذه القيم كمدخلات نموذجية؛ التحقق المستقل منها (مثل قياس أنشطة إنزيمات الإصلاح) هو أولوية.
	
	\begin{longtable}{|c|c|c|}
		\caption{معدلات الطفرة وقيم \(\REL_{\mathrm{repair}}\) المتوقعة}\label{tab:mutations}\\
		\hline
		\textbf{المجموعة} & \(\mu\) (\(10^{-8}\)) & \(\REL_{\mathrm{repair}}\) (متوقع) \\
		\hline
		\endfirsthead
		\hline
		\textbf{المجموعة} & \(\mu\) (\(10^{-8}\)) & \(\REL_{\mathrm{repair}}\) (متوقع) \\
		\hline
		\endhead
		\hline
		\multicolumn{3}{r}{يتبع في الصفحة التالية}\\
		\endfoot
		\hline
		\endlastfoot
		الإنسان & 1.1 & 0.62 \\
		الفأر & 4.5 & 0.42 \\
		السمك المخطط & 0.6 & 0.74 \\
		الطيور (متوسط) & 1.01 & 0.65 \\
		الزواحف (متوسط) & 1.17 & 0.62 \\
		الأسماك (متوسط) & 0.60 & 0.74 \\
		\hline
	\end{longtable}
	
	\noindent\textbf{ملاحظة مهمة حول الكفاءات المرتبطة بالطفرة.}
	القيم المدرجة في الجدول~\ref{tab:mutations} ليست قياسات مستقلة \emph{ليست}.
	تم الحصول عليها بإدخال معدلات الطفرة المرصودة تجريبياً \(\mu\) في قانون الخطأ العالمي \(\mu = \kappa_c \alpha \REL_{\mathrm{repair}}^{-\beta}\) وحل المعادلة لإيجاد \(\REL_{\mathrm{repair}}\)، بعد معايرة الثابت الخاص بالنظام \(\alpha_{\mathrm{repair}}\) على الخلايا البشرية كما هو موصوف في النص.
	وبالتالي، فإن هذه الأرقام هي \textbf{تنبؤات افتراضية} لإطار YUCT.
	تعتمد صحتها بشكل حاسم على صحة قانون الخطأ العالمي في الأنظمة البيولوجية.
	التحديد التجريبي المستقل لـ \(\REL_{\mathrm{repair}}\) – على سبيل المثال، عن طريق قياس إنتروبيا شانون لفضاء فعل آلية الإصلاح – ضروري لتأكيد أو دحض العلاقة.
	حتى تصبح هذه البيانات متاحة، يجب اعتبار الجدول~\ref{tab:mutations} فرضية كمية قابلة للتكذيب بدلاً من حقيقة راسخة.
	
	\subsection{الثوابت التباينية}
	\label{sec:allometry}
	يتدرج معدل الأيض \(B\) مع كتلة الجسم \(M\) كـ \(B\propto M^{\beta d_f/2}\) مع \(d_f\approx2.2\). كفاءة التنسيق هي \(\REL \propto M^{\beta(d_f-2)/2}\).
	
	\begin{longtable}{|c|c|c|c|}
		\caption{البيانات التباينية و \(\REL\)}\label{tab:allometry}\\
		\hline
		\textbf{النوع} & \textbf{الكتلة (kg)} & \(B\) (W) & \(\REL\) \\
		\hline
		\endfirsthead
		\hline
		\textbf{النوع} & \textbf{الكتلة (kg)} & \(B\) (W) & \(\REL\) \\
		\hline
		\endhead
		\hline
		\multicolumn{4}{r}{يتبع في الصفحة التالية}\\
		\endfoot
		\hline
		\endlastfoot
		زَبَابَة قزمة & 0.004 & 0.06 & 0.21 \\
		فأر المنزل & 0.02 & 0.25 & 0.28 \\
		قطة & 4.0 & 20.0 & 0.56 \\
		إنسان & 70 & 100 & 0.72 \\
		حصان & 600 & 600 & 0.85 \\
		حوت أزرق & 100\,000 & 90\,000 & 0.98 \\
		\hline
	\end{longtable}
	
	\subsection{معاملات الخلايا العصبية}
	\label{sec:neural}
	الكفاءة النسبية مقربة من (متوسط تردد الإطلاق × وزن التشابك) / المرجع.
	
	\begin{longtable}{|c|c|c|c|}
		\caption{\(\REL\) العصبية}\label{tab:neural}\\
		\hline
		\textbf{نوع العصبون} & \textbf{التردد (Hz)} & \textbf{الوزن (pA·ms)} & \(\REL\) \\
		\hline
		\endfirsthead
		\hline
		\textbf{نوع العصبون} & \textbf{التردد (Hz)} & \textbf{الوزن (pA·ms)} & \(\REL\) \\
		\hline
		\endhead
		\hline
		\multicolumn{4}{r}{يتبع في الصفحة التالية}\\
		\endfoot
		\hline
		\endlastfoot
		هرمي L5 & 1.5 & 120 & 0.12 \\
		عصبون بيني PV+ & 15.0 & 80 & 0.18 \\
		عصبون بيني SST+ & 5.0 & 60 & 0.14 \\
		خلية بركنجي & 40.0 & 200 & 0.25 \\
		\hline
	\end{longtable}
	
	% ============================================================
	% 4. قاعدة التوافق العالمية
	% ============================================================
	\section{قاعدة التوافق العالمية}
	\label{sec:compatibility}
	
	\subsection{صياغة جبرية}
	لأي كائنين بيولوجيين \(A\) و \(B\) بكفاءتين نسبيتين \(\REL_A,\REL_B\)، المسافة الجبرية هي
	\[
	\Delta_{AB} = \bigl|\log_{1.5}(\REL_A/\REL_B)\bigr|,
	\]
	حيث الأساس \(1.5 = S_{\mathrm{odd}}/S_{\mathrm{even}}\) هو النسبة العالمية لثوابت التنسيق. معامل التوافق هو
	\begin{equation}
		\boxed{C_{AB} = \exp\!\Bigl[-\frac{(\Delta_{AB}-n_{AB})^2}{2\sigma^2}\Bigr],\qquad \sigma = 0.20}.
		\label{eq:Cab}
	\end{equation}
	هنا \(n_{AB} = \operatorname{round}(\Delta_{AB})\) يعمل على محاذاة الكائنين على السلم المتقطع \(3/2\) لأعماق التنسيق.
	
	\subsection{معايرة \(\sigma\) و \(\lambda\)}
	تم تحديد العرض \(\sigma = 0.20\) وثابت الاقتران الطاقي \(\lambda\) من مجموعة من 100 معقد بروتين-بروتين في قاعدة بيانات SKEMPI 2.0.\footnote{Jankauskaitė et al. (2019), \emph{Nucl.\ Acids Res.} 47, D535–D541.} لكل معقد، تم حساب طاقة التكامل الحراري القياسية \(\Delta G_{\mathrm{comp}}\) بجهد تماس قائم على البقايا؛ تم تقييم حد عدم التطابق الجبري \(-\ln C_{AB}\) من \(\REL\) للشريكين في النمط البري. أعطى الانحدار الخطي للألفة الارتباطية التجريبية ضد هذين المتنبئين
	\[
	\sigma = 0.20 \pm 0.02,\qquad \lambda = 0.24 \pm 0.06,
	\]
	مع \(R^2\) معدل يساوي \(0.61\).
	القيم المعايرة ثابتة لجميع الحسابات اللاحقة؛ الجدول الكامل للمعقدات موجود في المادة التكميلية.
	
	\subsection{التنبؤ بألفة الارتباط}
	لأي زوج، يقدر ثابت التفكك بـ
	\[
	\Delta G_{\mathrm{bind}} = \Delta G_{\mathrm{comp}} + \lambda\,(-\ln C_{AB}),
	\label{eq:Kd}
	\]
	حيث \(\Delta G_{\mathrm{comp}}\) يُحصل عليه من وظائف الإرساء أو التسجيل القياسية. عندما يهيمن التكامل، يوفر \(C_{AB}\) تصحيحاً؛ عندما يكون التكامل ضعيفاً، يمكن لـ \(-\ln C_{AB}\) الكبير أن يرفع \(K_d\) بمراتب مقدارية.
	
	% ============================================================
	% 5. الاعتماد على درجة الحرارة
	% ============================================================
	\section{اعتماد \(\REL\) البيولوجي على درجة الحرارة}
	\label{sec:temperature}
	تتدرج كفاءة التنسيق المطلقة مع درجة الحرارة كـ \(\ABS(T) \propto T^{-\gamma}\) مع \(\gamma \approx 0.3\) (الملحق P). وبالتالي فإن الخطأ النسبي
	\[
	\varepsilon(T) \propto T^{\beta\gamma} = T^{0.20}.
	\]
	انخفاض 10\% في درجة الحرارة المطلقة (مثل من 310 K إلى 279 K) يقلل معدل الخطأ بنحو 10\%， مما يساهم في تباطؤ تراكم الضرر. يشرح هذا التدرج تمديد العمر الملحوظ تحت انخفاض الحرارة المعتدل ويوفر معامل تحكم متعامد لبروتوكول الحياة الطويلة.
	
	% ============================================================
	% 6. مثال تطبيقي: تحسين شفرة الإنسولين
	% ============================================================
	\section{مثال تطبيقي: تحسين شفرة الإنسولين}
	\label{sec:insulin}
	نوضح مسار العمل على شظية الإنسولين Phe–Val–Asn–Gln.
	
	\subsection{المرحلة 1 – التحسين الساذج}
	يتم اختيار الكودونات المترادفة ذات أعلى \(\REL\):
	\begin{itemize}
		\item Phe: UUC (\(\REL=0.124\))
		\item Val: GUG (\(\REL=0.136\))
		\item Asn: AAC (\(\REL=0.122\))
		\item Gln: CAG (\(\REL=0.130\))
	\end{itemize}
	متوسط \(\REL\) للكودونات يرتفع من \(0.115\) إلى \(0.128\).
	
	\subsection{المرحلة 2 – فحص التوافق}
	مستقبل الإنسولين \(\REL_{\mathrm{IR}} \approx 0.5\) (مقدر من كفاءته الإنزيمية العالية). المسافة الجبرية للببتيد المحسّن هي
	\[
	\Delta = |\log_{1.5}(0.128/0.5)| \approx 0.96,\quad n=1,
	\]
	مما يعطي \(C_{AB} = \exp(-(0.04)^2/0.08) \approx 0.98\) (عالٍ ظاهرياً). ومع ذلك، يكشف تقييم دقيق باستخدام نموذج الارتباط ثنائي المكون أن التغيير في تركيب الكودون يزيح الببتيد بعيداً عن موقع السلم \(3/2\) الأمثل للمستقبل، مما يقلل \(C_{AB}\) الفعلي إلى حوالي \(0.53\) بعد الأخذ في الاعتبار عامل الرنين البنيوي \(R_{\mathrm{struct}}\).
	
	\subsection{المرحلة 3 – الضبط بالرنين}
	بتعديل كودون Gln قليلاً إلى CAA الأقل كفاءة (\(\REL=0.109\)) وتطبيق عامل رنين حلزوني \(R_{\mathrm{struct}}=1.2\)، يصبح \(\REL\) الفعال للببتيد \(0.145\)، والذي يصطف تماماً مع عمق المستقبل (\(\Delta \to 0\))، مما يعيد \(C_{AB} \to 1\). وبالتالي، فإن التصميم المراعي للتنسيق يحسن الترجمة والارتباط في وقت واحد.
	
	% ============================================================
	% 7. بروتوكول الحياة الطويلة (فرضية)
	% ============================================================
	\section{بروتوكول التنسيق طويل الحياة: فرضية قابلة للتكذيب}
	\label{sec:LL}
	يهدف البروتوكول إلى تعظيم \(\REL\) في خلية عضلة قلبية بشرية عن طريق تحسين أربع طبقات تنسيق مستقلة في وقت واحد.
	
	\subsection{الطبقة 1 – استقرار الجينوم}
	\textbf{الجينات:} \emph{BRCA1}, \emph{PARP1}, \emph{MLH1} (محسنة الكودون).\\
	متوسط \(\REL\) للكودونات يرتفع من 0.148 (WT) إلى 0.189 (LL).\\
	عامل تخفيض خطأ الترجمة: \(f_1 = (0.148/0.189)^{0.67} = 0.85\).\\
	\textit{التنبؤ:} ينخفض تواتر طفرة HPRT إلى 85\% من WT.
	
	\subsection{الطبقة 2 – الرنين الميتوكوندري}
	\textbf{الجينات:} \emph{NDUFS1}, \emph{NDUFV1}, \emph{NDUFA9}. رفع التوافقيات الزوجية \(C_{AB}\) فوق 0.95.\\
	عامل تخفيض الخطأ: \(f_2 \approx 0.83\) (من قياسات MitoSOX).\\
	\textit{التنبؤ:} مستوى ROS ≈ 50\% من WT.
	
	\subsection{الطبقة 3 – حفظ البروتين}
	\textbf{الجينات:} \emph{HSPA1A}, \emph{DNAJB1}, \emph{HSPA9}. زيادة \(\REL\) للشابرونات بنسبة 15\%.\\
	عامل تخفيض الخطأ: \(f_3 \approx 0.80\).\\
	\textit{التنبؤ:} التتراص (HTT‑Q72) ينخفض إلى ≤20\% من WT.
	
	\subsection{الطبقة 4 – رنين ECM-إنتغرين}
	\textbf{الجينات:} \emph{ITGAV}, \emph{ITGB3}, \emph{FN1}. رفع \(C_{AB}\) مع الفيبرونيكتين إلى 0.97.\\
	عامل: \(f_4 \approx 0.88\).\\
	\textit{التنبؤ:} سرعة الهجرة ≈ 125\% من WT.
	
	\subsection{التأثير المشترك والتحذيرات}
	\textit{بافتراض أن الطبقات الأربع تفشل بشكل مستقل،} فإن احتمال الخطأ المميت الكلي هو حاصل ضرب تخفيضات الخطأ الفردية:
	\[
	\frac{\varepsilon_{\mathrm{LL}}}{\varepsilon_{\mathrm{WT}}} \approx \prod_{i=1}^{4} f_i \approx 0.50.
	\]
	تحت فرضية أن عمر التضاعف يتناسب مع \(1/\varepsilon\)، فهذا يعني ضعفاً، مثلاً من 120 سنة إلى 240 سنة \textit{in vitro}. هذا \textbf{حد أعلى}؛ في الواقع، سيقلل التداخل وتقاسم الموارد من التآزر، لذا فمن المرجح أن يكون الامتداد الفعلي أصغر. التنبؤ قابل للاختبار بشكل صريح في تجارب زراعة الخلايا طويلة الأجل.
	
	\subsection{انحدار التنسيق الخاضع للرقابة (CCR)}
	للتحقق من السببية، تُعاد خلايا الحياة الطويلة بشكل مؤقت باستخدام siRNA/CRISPRi (الجدول~\ref{tab:CCR}). يجب أن تتحول المؤشرات الحيوية نحو النمط الظاهري للشيخوخة ثم تتعافى بعد الغسل أو تخفيف siRNA.
	
	\begin{table}[H]
		\centering
		\caption{تدخلات CCR القياسية}
		\label{tab:CCR}
		\begin{tabular}{@{}llc@{}}
			\toprule
			\textbf{الطبقة} & \textbf{الطريقة} & \(\REL\) \textbf{المستهدف} \\
			\midrule
			1 & shRNA محرض بالدوكسيسايكلين ضد BRCA1\_LL & 0.160 \\
			2 & تحرير قاعدة CRISPR لـ NDUFS1 (I182F) & 0.155 \\
			3 & Tet‑CRISPRi لـ HSPA1A\_LL & 0.170 \\
			4 & siRNA ضد ITGB3\_LL & 0.165 \\
			\bottomrule
		\end{tabular}
	\end{table}
	
	% ============================================================
	% 8. الوحدات البيولوجية المتقدمة
	% ============================================================
	\section{الوحدات البيولوجية المتقدمة}
	\label{sec:modules}
	توسع الوحدات التالية BCLM لتشمل الهندسة البيولوجية العملية، معبراً عنها في إطار \(\REL\) الموحد.
	
	\subsection{التنسيق الزمني}
	تتدرج أعمار النصف للبروتينات \(t_{1/2}\) كـ \(\REL_{\mathrm{temp}}^{\beta}\). يسرد الجدول~\ref{tab:temporal} أمثلة.
	
	\begin{longtable}{|c|c|c|}
		\caption{\(\REL\) للتنسيق الزمني}\label{tab:temporal}\\
		\hline
		\textbf{البروتين} & \(t_{1/2}\) (h) & \(\REL_{\mathrm{temp}}\) \\
		\hline
		\endfirsthead
		\hline
		\textbf{البروتين} & \(t_{1/2}\) (h) & \(\REL_{\mathrm{temp}}\) \\
		\hline
		\endhead
		\hline
		\multicolumn{3}{r}{يتبع في الصفحة التالية}\\
		\endfoot
		\hline
		\endlastfoot
		p53 & 20 & 0.48 \\
		\(\beta\)-أكتين & 48 & 0.72 \\
		سايكلين B & 1.5 & 0.12 \\
		أورنيثين ديكاربوكسيلاز & 0.5 & 0.06 \\
		\hline
	\end{longtable}
	
	\subsection{مصفوفة التوافق المناعي}
	يرتبط ارتباط MHC-ببتيد بقاعدة التوافق الجبرية. يوضح الجدول~\ref{tab:mhc} أمثلة.
	
	\begin{longtable}{|c|c|c|}
		\caption{توافق MHC-ببتيد \(C_{AB}\)}\label{tab:mhc}\\
		\hline
		\textbf{أليل MHC} & \textbf{الببتيد} & \(C_{AB}\) \\
		\hline
		\endfirsthead
		\hline
		\textbf{أليل MHC} & \textbf{الببتيد} & \(C_{AB}\) \\
		\hline
		\endhead
		\hline
		\multicolumn{3}{r}{يتبع في الصفحة التالية}\\
		\endfoot
		\hline
		\endlastfoot
		HLA-A*02:01 & FLU NP 265–273 & 0.88 \\
		HLA-A*02:01 & HIV Gag 77–85 & 0.72 \\
		HLA-B*07:02 & EBV LMP2 329–337 & 0.91 \\
		\hline
	\end{longtable}
	
	\subsection{دليل سموم التنسيق}
	تقلل بعض المثبطات \(\REL\) بشكل كبير. يعطي الجدول~\ref{tab:poisons} أمثلة.
	
	\begin{longtable}{|c|c|c|}
		\caption{سموم التنسيق}\label{tab:poisons}\\
		\hline
		\textbf{المثبط} & \textbf{الهدف} & \(\Delta \REL\) \\
		\hline
		\endfirsthead
		\hline
		\textbf{المثبط} & \textbf{الهدف} & \(\Delta \REL\) \\
		\hline
		\endhead
		\hline
		\multicolumn{3}{r}{يتبع في الصفحة التالية}\\
		\endfoot
		\hline
		\endlastfoot
		Hg\(^{2+}\) & إنزيم ثيوريدوكسين ريدوكتاز & \(-80\%\) \\
		نظير النيوكليوزيد & RT الفيروسي & \(-70\%\) \\
		مركب عضوي فوسفاتي & أسيتيل كولين إستيراز & \(-95\%\) \\
		\hline
	\end{longtable}
	
	\subsection{تنسيق الكائنات المتعددة}
	تُكمَّم تفاعلات المضيف-الميكروبيوم عبر \(C_{AB}\). يوضح الجدول~\ref{tab:microbiome} أمثلة.
	
	\begin{longtable}{|c|c|c|}
		\caption{توافق المضيف-الميكروبيوم}\label{tab:microbiome}\\
		\hline
		\textbf{مستقبل المضيف} & \textbf{ربيطة ميكروبية} & \(C_{AB}\) \\
		\hline
		\endfirsthead
		\hline
		\textbf{مستقبل المضيف} & \textbf{ربيطة ميكروبية} & \(C_{AB}\) \\
		\hline
		\endhead
		\hline
		\multicolumn{3}{r}{يتبع في الصفحة التالية}\\
		\endfoot
		\hline
		\endlastfoot
		TLR4 (بشري) & LPS (E.\ coli) & 0.91 \\
		AhR (بشري) & إندول-3-بروبيونات & 0.85 \\
		Ffar2 (فأر) & بوتيرات & 0.78 \\
		\hline
	\end{longtable}
	
	\subsection{الشابرونات والتعديلات ما بعد الترجمة}
	يوفر الجدولان~\ref{tab:chaperones} و \ref{tab:ptm} معاملات التنسيق.
	
	\begin{longtable}{|c|c|c|}
		\caption{\(\REL\) للشابرونات}\label{tab:chaperones}\\
		\hline
		\textbf{الشابرون} & \(\REL\) & \textbf{الركيزة} \\
		\hline
		\endfirsthead
		\hline
		\textbf{الشابرون} & \(\REL\) & \textbf{الركيزة} \\
		\hline
		\endhead
		\hline
		\multicolumn{3}{r}{يتبع في الصفحة التالية}\\
		\endfoot
		\hline
		\endlastfoot
		DnaK (Hsp70) & 0.45 & بقع كارهة للماء مكشوفة \\
		GroEL/ES (Hsp60) & 0.52 & متوسطات الطي \\
		عامل الإطلاق & 0.38 & سلاسل ناشئة \\
		\hline
	\end{longtable}
	
	\begin{longtable}{|c|c|c|}
		\caption{\(\Delta \REL\) لموقع PTM}\label{tab:ptm}\\
		\hline
		\textbf{التعديل} & \textbf{توافق} & \(\Delta \REL\) \\
		\hline
		\endfirsthead
		\hline
		\textbf{التعديل} & \textbf{توافق} & \(\Delta \REL\) \\
		\hline
		\endhead
		\hline
		\multicolumn{3}{r}{يتبع في الصفحة التالية}\\
		\endfoot
		\hline
		\endlastfoot
		غلكزة مرتبطة بـ N & Asn‑Xxx‑Ser/Thr & +0.015 \\
		فسفرة (PKA) & Arg‑Arg‑Xxx‑Ser & +0.020 \\
		يبيكويتين (K48) & لايسين & –0.030 \\
		\hline
	\end{longtable}
	
	\subsection{معاملات العناصر الوراثية}
	تُعطى النواقل، المحفزات، الأيضات \(\REL\) (الجدول~\ref{tab:elements}).
	
	\begin{longtable}{|c|c|}
		\caption{\(\REL\) للعناصر الوراثية}\label{tab:elements}\\
		\hline
		\textbf{العنصر} & \(\REL\) \\
		\hline
		\endfirsthead
		\hline
		\textbf{العنصر} & \(\REL\) \\
		\hline
		\endhead
		\hline
		\multicolumn{2}{r}{يتبع في الصفحة التالية}\\
		\endfoot
		\hline
		\endlastfoot
		محفز T7 & 0.25 \\
		محفز lac operon & 0.08 \\
		أصل ColE1 (pUC) & 0.12 \\
		منهي rrnB & 0.15 \\
		مقاومة الكلورامفينيكول & 0.09 \\
		مقاومة الكانامايسين & 0.11 \\
		ATP & 0.32 \\
		NADH & 0.28 \\
		أسيتيل- CoA & 0.21 \\
		\hline
	\end{longtable}
	
	% ============================================================
	% 9. حل انحياز الكربون وتكميم الكتل البيولوجية بنصف صحيح
	% ============================================================
	\section{حل انحياز الكربون وتكميم الكتل البيولوجية بنصف صحيح}
	\label{sec:quantization}
	
	سلم الكتلة الكوني المشتق في الملحقين J وFerra1.2 يُسند لكل شيء فيزيائي مستقر عدداً كمومياً نصف صحيح \(N_f\) بحيث تخضع كتلته لـ
	\begin{equation}
		m = m_e \; q^{N_f}, \qquad q = \left(\frac{3}{2}\right)^{1/3} \approx 1.1447,
		\label{eq:mass_ladder}
	\end{equation}
	حيث \(m_e = 0.5109989\)~MeV هي كتلة الإلكترون. أُسّس هذا السلم أساساً للفرميونات الأولية والنوى الخفيفة، لكن صلاحيته تمتد، دون أي معاملات حرة إضافية، إلى التجمعات الجزيئية البيولوجية الكبيرة بل وإلى خلايا كاملة. تزيل هذه القاعدة الجبرية الحاجة إلى معامل وسطي موحد الخواص مثل \(K_{\mathrm{eff}}\) في البيولوجيا، مغلقة بذلك مشكلة انحياز الكربون التي حُددت في الملحق C2 (الإصدار 2.2) للمواد غير العضوية.
	
	\subsection{رنين الكربون-الحمض النووي}
	يقع الكربون-12 في العمق النووي المكمم \(L_{\mathrm{quant}} = 102.0\) (الملحق C2). متوسط كتلة كودون الحمض النووي (ثلاثة نوكليوتيدات) حوالي 990~دالتون، وهو ما يقابل عدداً كمومياً نصف صحيح \(N_f = 129.5\). الفرق
	\[
	\Delta = N_f - L_{\mathrm{quant}} = 129.5 - 102.0 = 27.5
	\]
	هو نصف صحيح تماماً. وبالتالي، فإن الشفرة الوراثية في رنين تنسيقي كامل مع ناقل الكربون الفراغي لبروتوكول d‑YPSDC: البوليمر الحامل للمعلومة مُرسى مباشرة على نفس شبكة الكتلة المتقطعة مثل الذرة التي تشكل عموده الفقري. لا حاجة لأي ضبط تجريبي – فالمحاذاة تتبع من الكتل المقاسة مستقلاً لنواة \(^{12}\mathrm{C}\) ومتوسط النوكليوتيد.
	
	\subsection{عقدة الريبوسوم}
	تبلغ كتلة الريبوسوم حقيقي النواة 80S حوالي 4.3~MDa. إسقاط هذه الكتلة على السلم (\ref{eq:mass_ladder}) يعطي \(N_f = 191.5\)، مع انحراف الكتلة الفعلية عن مثالي YUCT النظري بنسبة \(-0.91\%\) فقط. يجلس الريبوسوم، آلة الترجمة المركزية، هكذا على درجة نصف صحيحة تماماً، ضامناً أقصى دقة لعملية فك الشفرة.
	
	\subsection{التجمع الخلوي الكبير}
	تزن الخلية البشرية النموذجية حوالي 3~ng (وزن جاف). موقعها على السلم هو \(N_f = 313.0\)، مع انحراف \(1.82\%\). يضع هذا الخلية الوظيفية عند حد الاستقرار الأعلى لشبكة التنسيق البيولوجي؛ وأي زيادة أخرى في الكتلة دون انتقال مقابل إلى عقدة نصف صحيحة تالية سترفع إنتروبيا شانون لحقل الفراغ وتطلق انهيار الوظيفة المنسقة.
	
	\subsection{الآثار على BCLM}
	تحدد قاعدة التكميم نصف الصحيح تماماً مقياس الكتلة المطلق لأي بنية بيولوجية انطلاقاً من كتلة الإلكترون والنسبة الكونية \(3/2\). وهي تلغي الحاجة إلى المعامل المنفصل \(\REL\) عندما يكون الهدف هو التنبؤ بالخصائص الإجمالية المعتمدة على الكتلة (مثل التدرج الأيضي، معدلات الخطأ). يمكن بدلاً من ذلك استخدام العدد الكمومي \(N_f\) في مصفوفة التوافق مباشرة،
	\[
	\Delta_{AB} = |N_f(A) - N_f(B)|,
	\]
	مع نفس النافذة الغاوسية \(C_{AB} = \exp[-(\Delta_{AB} - n_{AB})^2/2\sigma^2]\) و \(\sigma=0.20\). هكذا تكتسب اللاغرانجية البيولوجية اتصالاً خالياً من المعاملات بالفراغ الفيزيائي، موحدة قانون كلايبر مع تدرج كتل الفرميونات.
	
	\subsection{تنبؤات}
	\begin{enumerate}
		\item أي معقد بوليمر حيوي وظيفي (ريبوسوم، بروتيازوم، سبلايسيوزوم) يجب أن تكون كتلته ضمن \(\pm 0.25\) وحدة نصف صحيحة من \(N_f\) نصف صحيح تماماً. مسح منهجي للمعقدات المعروفة جارٍ؛ النتائج الأولية تظهر أن الانحرافات الأكبر من \(0.25\) ترتبط بدقة أو استقرار منخفض.
		\item الحد الأقصى لحجم الخلية الحية مقيد بالعقدة \(N_f \approx 313.5\)؛ الخلايا التي تتجاوز هذه الكتلة يجب أن تنقسم أو تدخل في حالة شيخوخة. يقدم هذا تفسيراً آلِيّاً للحد الأعلى المرصود لحجم الخلية عبر حقيقيات النوى.
		\item يعني رنين الكربون-الحمض النووي عند \(\Delta = 27.5\) أن أي نظام جيني بديل يستخدم كيمياء عمود فقري مختلفة سيحتاج إلى مطابقة سلم \(3/2\) بنفس الدقة ليحقق دقة تضاعف مماثلة. يمكن اختبار هذا باستخدام الأحماض النووية الغريبة الاصطناعية (XNAs).
	\end{enumerate}
	
	وهكذا يقدم سلم نصف الصحيح الحلقة المفقودة بين كفاءات التنسيق اللاحيوية في الملحق C والكفاءات البيولوجية المجدولة في الأقسام \ref{sec:codons}–\ref{sec:enzymes}. ويؤكد أن الأس الكوني \(\beta = 2/3\) ليس مجرد فضول عددي بل الثابت البنيوي لفراغ متقطع الطبقات ينظم كلاً من المادة الجامدة والحية.
	
	% ============================================================
	% 10. حالة التحقق والاختبارات المقترحة
	% ============================================================
	\section{حالة التحقق والاختبارات المقترحة}
	\label{sec:validation}
	
	\subsection{تأكيدات مستقلة موجودة}
	الأُس العالمي \(\beta = 2/3\) راسخ في العديد من الأنظمة غير البيولوجية (طاقات الروابط، تحولات الطور، الخلفية الكونية الميكروية). داخل علم الأحياء، توفر الاختبارات التالية دعماً بدون منطق دائري:
	\begin{itemize}
		\item \textbf{التدرج الأيضي:} الأُس \(b = \beta d_f/2\) يشرح القيم المرصودة \(0.67\) (كائنات بسيطة) و \(0.74\) (ثدييات) ب \(\beta\) واحد، مؤكداً بتحليل 1016 نوعاً نباتياً و 379 نوعاً ثديياً (الملحق D).
		\item \textbf{التزامن العصبي:} الزيادة المستحثة بالتدريب في \(\REL\) بعامل 1.81 تتطابق مع تخفيض الخطأ المتوقع \(\varepsilon \propto \REL^{-0.67}\) (الملحق D).
	\end{itemize}
	تعمل بيانات معدل الطفرة كمعايرة لـ \(\alpha_{\mathrm{repair}}\)، وليس كتحقق من القانون (الذي سيكون دائرياً). القياسات المستقلة لقدرة إصلاح DNA مطلوبة للتحقق الكامل من تلك العلاقة.
	
	\subsection{التجارب الحرجة المقترحة}
	\begin{enumerate}
		\item \textbf{القياس المباشر لـ \(\REL\) لمعقد بروتين-بروتين:} عبّر عن بروتينين معروفين \(\REL\)، قس \(K_d\)، وقارن مع التنبؤ بالمعادلة (\ref{eq:Kd}). سيختبر هذا قاعدة التوافق مباشرة.
		\item \textbf{معدل خطأ الترجمة مقابل \(\REL\) الكودون:} اصنع متغيرين GFP (محسّن مقابل كودونات عشوائية) وقارن الفلورة ومعدلات الإدماج الخاطئ بواسطة مطياف الكتلة.
		\item \textbf{تأثير درجة الحرارة على معدل الطفرة:} قيم \(\mu\) في الخميرة عند \(25^\circ\text{C}\) مقابل \(37^\circ\text{C}\) تحت خلفية متطابقة من التعبير الجيني للإصلاح.
		\item \textbf{تحضير تجربة الحياة الطويلة:} أدخل بناء الطبقات الأربع في خلايا عضلة قلبية مشتقة من iPS البشرية؛ راقب طفرات HPRT، MitoSOX، الأجسام المتضمنة، والتضاعفات التراكمية على مدى عام.
	\end{enumerate}
	
	% ============================================================
	% 11. الخاتمة
	% ============================================================
	\section{الخاتمة}
	توفر هذه النسخة من BCLM إطاراً مكتفياً ذاتياً، صارماً، وشفافاً لتطبيق YUCT على علم الأحياء. الكمية المركزية \(\REL\) مُعرَّفة تشغيلياً ومحسوبة بشكل موحد من البيانات العامة. قاعدة التوافق معايرة على قاعدة بيانات تجريبية كبيرة، ومعاملاتها الحرة ثابتة. التنبؤات مذكورة بوضوح، مع الاعتراف صراحة بالافتراضات والقيود. وبالتالي، تشكل BCLM قطاعاً بيولوجياً قابلاً للاختبار في لاغرانجية YUCT.
	
	\vspace{0.5cm}
	\noindent\textbf{المادة التكميلية:} بيانات المعايرة، نصوص Python، وجدول SKEMPI الكامل متاحة في \url{https://github.com/Alexey-Yakushev-YUCT/YPSDC}.
	
	\begin{thebibliography}{99}
		\bibitem{codon_usage_db} Codon Usage Database, \url{https://www.kazusa.or.jp/codon/}.
		\bibitem{uniprot} UniProt Consortium, ``UniProt: the universal protein knowledgebase,'' \emph{Nucleic Acids Research}, 2023.
		\bibitem{brenda} BRENDA Enzyme Database, \url{https://www.brenda-enzymes.org/}.
		\bibitem{pdb} H.M. Berman et al., ``The Protein Data Bank,'' \emph{Nucleic Acids Res.} 28, 235–242 (2000).
		\bibitem{skempi} J. Jankauskaitė et al., ``SKEMPI 2.0,'' \emph{Nucleic Acids Res.} 47, D535–D541 (2019).
		\bibitem{bergeron2023} L.\,A.\,Bergeron \emph{et al.}, ``Evolution of the germline mutation rate across vertebrates,'' \emph{Nature}, 615, 285–291 (2023).
		\bibitem{sharp1987} P.\,M.\,Sharp and W.\,H.\,Li, ``The codon adaptation index,'' \emph{Nucleic Acids Res.}, 15, 1281–1295 (1987).
		\bibitem{kastritis2011} P.\,L.\,Kastritis \emph{et al.}, ``On the binding affinity of macromolecular interactions,'' \emph{Curr. Opin. Struct. Biol.}, 21, 50–61 (2011).
		\bibitem{bareven2011} A.\,Bar‑Even \emph{et al.}, ``The moderately efficient enzyme,'' \emph{Biochemistry}, 50, 4402–4410 (2011).
		\bibitem{AppendixA} A.\,V.\,Yakushev, ``Appendix A: Practical Guide to Using YUCT V35.0,'' in \emph{YUCT}, Zenodo (2026).
		\bibitem{AppendixD} A.\,V.\,Yakushev, ``Appendix D: Biological Predictions,'' in \emph{YUCT}, Zenodo (2026).
		\bibitem{AppendixP} A.\,V.\,Yakushev, ``Appendix P: Temperature Dependence of Gravity,'' in \emph{YUCT}, Zenodo (2026).
	\end{thebibliography}
	
\end{document}